% Round 30 — Chương 1; tệp mới để ghép, không sửa luận văn gốc.
% Bằng chứng: EVIDENCE-REGISTER, ADDENDUM-1…4; C2-xx ở ADDENDUM-5-CH2.
% Không cần gói mới; TikZ, booktabs và khungy đã có trong main.tex.
\chapter{Đặt vấn đề}
\label{chuong:modau}

Chương này nêu vì sao giờ trả đã hứa với khách đi chung xe có thể bị dời, chia độ lệch đó thành ba phần và chỉ ra
chỗ khóa luận xử lý từng phần. Sau đó, chương nêu các câu hỏi nghiên cứu, đóng góp và phạm vi của khóa luận.
% EV: — (lộ trình chương).

\section{Bối cảnh: lời hứa giờ trả trong dịch vụ đi chung xe}
\label{muc:c1-boicanh}

Trong dịch vụ đi chung xe (ride-pooling), bộ điều phối ghép các yêu cầu vào cùng một xe và cập nhật kế hoạch khi
có yêu cầu mới. Giờ trả đã báo cho khách vì thế có thể đổi khi xe nhận thêm điểm dừng; ngay cả khi giữ nguyên
kế hoạch, thời gian đi thay đổi cũng có thể làm mốc giờ trả khác đi. Cần phân biệt hai phần này trước khi đặt
giới hạn cho việc sửa lời hứa.
% EV: J17; A01; A02; J01; J06.

Hai khách trong bản ghi mô phỏng ở Hình~\ref{hinh:c1-loihua} minh họa sự phân biệt đó. Khách ở hình trên trễ
307,3 giây so với lời hứa đầu: 99,1 giây là phần thay đổi khi giữ kế hoạch và 208,2 giây được tính cho điều phối.
Đây là khách trễ nhất trong một kịch bản phát triển với trần 300 giây. Khi dùng trần 30 giây, khách trễ nhất
của kịch bản là một người khác, trễ 99,1 giây, trong đó 97,1 giây thuộc phần giữ kế hoạch.
% EV: D06; N-30; W07:48-54; q4-rider-2026-09-30.txt (khách cực đại của hai cấu hình khác nhau).

\begin{figure}[htbp]
\centering
\begin{tikzpicture}[x=1cm,y=1cm,>=Stealth,font=\footnotesize]
  \node[anchor=west] at (0,4.85) {(a) Khách trễ nhất, trần 300 giây};
  \draw[->] (0.6,3.75)--(11.8,3.75);
  \draw (0.6,3.64)--(0.6,3.86);
  \draw (3.673,3.64)--(3.673,3.86);
  \draw (10.127,3.64)--(10.127,3.86);
  \fill (0.6,3.75) circle (1.8pt);
  \fill (3.673,3.75) circle (1.8pt);
  \fill (10.127,3.75) circle (1.8pt);
  \draw[<->] (0.6,4.18)--(3.673,4.18);
  \node[align=center] at (2.137,4.48) {Giữ kế hoạch\\99,1 giây};
  \draw[<->] (3.673,4.18)--(10.127,4.18);
  \node[align=center] at (6.9,4.48) {Phần tính cho điều phối\\208,2 giây};
  \node[anchor=north,align=center] at (0.6,3.55) {Lời hứa đầu\\0};
  \node[anchor=north,align=center] at (3.673,3.55) {Mốc giữ kế hoạch\\99,1};
  \node[anchor=north,align=center] at (10.127,3.55) {Trả thực tế\\307,3};
  \node[anchor=west] at (0,2.45) {(b) Khách có lời hứa tính lại theo hồ sơ lịch sử};
  \draw[->] (0.6,1.35)--(11.8,1.35);
  \draw (0.6,1.24)--(0.6,1.46);
  \draw (8.553,1.24)--(8.553,1.46);
  \draw (9.363,1.24)--(9.363,1.46);
  \fill (0.6,1.35) circle (1.8pt);
  \draw[fill=white] (8.553,1.35) circle (2.4pt);
  \fill (9.363,1.35) circle (1.8pt);
  \draw[<->] (0.6,1.8)--(8.553,1.8);
  \node at (4.58,2.04) {Lời hứa dài thêm 172,9 giây};
  \node[anchor=north,align=center] at (0.6,1.14) {Lời hứa đầu\\7:34:51\\0};
  \draw[dashed] (8.553,1.25)--(7.9,0.53);
  \node[anchor=north,align=center] at (7.9,0.48) {Lời hứa có dự báo\\7:37:44\\172,9};
  \draw[dashed] (9.363,1.25)--(10.45,0.94);
  \node[anchor=north,align=center] at (10.45,0.89) {Trả thực tế\\7:38:02\\190,5};
  \node[anchor=west,align=left] at (10.15,1.87) {Còn trễ\\17,6 giây};
\end{tikzpicture}
\caption{Hai khách trong bản ghi mô phỏng. Mỗi trục tính số giây sau lời hứa đầu, với tỉ lệ riêng.
Hình trên thuộc kịch bản phát triển ngày 13/11/2018, khung 7:00--9:00; mốc giữ kế hoạch được tính lại trong mô hình.
Hình dưới thuộc ngày 15/11/2018, cùng khung giờ; giờ đồng hồ làm tròn tới giây, độ lệch tính trước khi làm tròn.}
\label{hinh:c1-loihua}
\end{figure}
% EV: D06; N-30; C31; W07:48-54; EXAMPLE:1-4,60-61.

Khách ở hình dưới đặt xe lúc 7:11:56, được đón lúc 7:22:30 và tới nơi muộn hơn lời hứa đầu 190,5 giây.
Khi tính lại lời hứa từ hồ sơ tắc đường lịch sử, độ lệch còn 17,6 giây. Ví dụ này minh họa vai trò của dự báo,
không phải phép đo tác dụng của trần.
% EV: C31; EXAMPLE:1-4,60-61; REVIEW-DECISIONS-R20 LR-15.

Vấn đề giờ đến cũng xuất hiện trong nguồn hướng dẫn dịch vụ: UberX Share nêu rằng không bảo đảm giờ đến,
còn MOIA Express giới thiệu cửa sổ giờ đến ngắn hơn như một lựa chọn nâng cấp. Về phía người đi xe, khảo sát của
Alonso-González và cộng sự~\cite{alonsogonzalez2020} đo giá trị của độ tin cậy trong dịch vụ đi chung.
Phân tích tổng hợp của Wardman và cộng sự~\cite{wardman2026} đo giá trị của thời gian tới muộn theo những cách hỏi khác.
Hai loại bằng chứng này không cho biết khách sẽ coi trọng trần điều phối của khóa luận đến mức nào.
% EV: J28; J31; J24; J23; C2-18; C2-17; I07; I10.

\section{Ba phần của độ lệch giờ trả và chỗ xử lý mỗi phần}
\label{muc:c1-cauhoi}

Trên tập kiểm tra chính, khi điều phối không trần, 134 trong 203 khách trễ quá một phút từng có giờ trả bị
điều phối dời. Đây là giao của hai tập khách, không phải số người trễ được xác định là do điều phối.
Theo sổ, phần tính cho điều phối bằng 73,4\% tổng số giây trễ của cả 203 khách: 23.027 trên 31.353 giây,
gộp theo khách và làm tròn hai tổng tới giây. Một khách bị dời tổng cộng 590,1 giây, gần mười phút.
% EV: C21; C24; C4A-03 (23.027,239 / 31.353,407 s); C20.

Khóa luận chia độ lệch so với lời hứa đầu thành ba phần ở Hình~\ref{hinh:c1-baphan}. Phần tính cho điều phối
được đo bằng cách so kế hoạch được chọn với kế hoạch cũ tính lại ở cùng trạng thái. Luật C giới hạn tổng độ lớn
của các lần dời thuộc phần này cho từng khách. Phần hồ sơ tắc đường lịch sử giải thích được cộng sẵn vào lời hứa
theo dự báo. Phần còn lại được ghi trong sổ theo từng khách để biết dự báo đang dùng chưa giải thích được bao nhiêu.
% EV: A01; A02; A06; F01; F09; X3-12; muc:so của Chương 3 v3.

\begin{figure}[htbp]
\centering
\begin{tikzpicture}[>=Stealth,font=\small,
  part/.style={draw,rounded corners,align=center,text width=5.6cm,minimum height=0.85cm},
  place/.style={draw,align=center,text width=5.2cm,minimum height=0.85cm}]
  \node[part] (dispatch) at (0,2.3) {Phần tính cho điều phối};
  \node[part] (history) at (0,1.15) {Phần hồ sơ lịch sử giải thích};
  \node[part] (rest) at (0,0) {Phần còn lại};
  \node[place] (cap) at (7.4,2.3) {Luật C: giới hạn từng khách};
  \node[place] (forecast) at (7.4,1.15) {Lời hứa theo dự báo};
  \node[place] (ledger) at (7.4,0) {Sổ theo từng khách: ghi và kiểm};
  \draw[->] (dispatch)--(cap);
  \draw[->] (history)--(forecast);
  \draw[->] (rest)--(ledger);
\end{tikzpicture}
\caption{Ba phần của độ lệch giờ trả và chỗ xử lý mỗi phần. Sổ ghi cả ba phần; các mũi tên chỉ vai trò
được nhấn mạnh, không biểu diễn quan hệ nhân quả.}
\label{hinh:c1-baphan}
\end{figure}
% EV: A06; F01; F09; X3-12; eq:so của Chương 3 v3.

Phép chia này là phép tính theo mô hình. Phần còn lại phụ thuộc hồ sơ dự báo đang dùng; nó không đồng nghĩa với
phần không thể biết trước. Từ cách chia trên, khóa luận đặt ba câu hỏi:
% EV: F09; K04; K19; X3-12.
\begin{description}[leftmargin=1.15cm,style=nextline,itemsep=3pt]
\item[CH1.] Phần dời do điều phối được đo trên đại lượng nào, kiểm khi nào, và trần trên phần đó giữ được gì,
không giữ được gì? Mục~\ref{muc:luatc} nêu cách đặt luật; các mục~\ref{muc:c4-dieuphoi},
\ref{muc:c4-congdon} và~\ref{muc:c4-hople} kiểm nó. Câu hỏi này không suy ra bảo đảm giờ đến thực tế.
\item[CH2.] Trên các tập đã đo, giữ trần đổi số lượt phục vụ, số lượt đúng giờ và độ dài lời hứa thế nào so với
điều phối không trần, luật hạn chót và cách đệm lời hứa? Các mục~\ref{muc:c4-ketqua}, \ref{muc:c4-chontran},
\ref{muc:c4-vaitro} và~\ref{muc:c4-dem} trả lời; kết quả không xác định tác dụng riêng của bộ chọn hay mức trần tốt nhất.
\item[CH3.] Lời hứa theo hồ sơ tắc đường lịch sử thay đổi số lượt đúng giờ thế nào, và phần nào còn chưa được
giải thích? Các mục~\ref{muc:dubao}, \ref{muc:so}, \ref{muc:c4-dubao} và~\ref{muc:c4-tuyet} trả lời;
kết quả không suy ra phản ứng của khách.
\end{description}
% EV: OUTLINE-V2 cuối §1.2 (câu hỏi và lộ trình); K05; K19; H02; I07.

\section{Các công trình gần nhất và chỗ đứng của khóa luận}
\label{muc:c1-lienquan}

Các nghiên cứu trước đã đo độ bất ổn, kiểm soát việc đổi giờ trước quyết định và đặt giới hạn theo từng khách.
Bảng~\ref{bang:c1-chodung} đối chiếu đại lượng được xét và cách xử lý phần thay đổi khi giữ kế hoạch.
Chương~\ref{chuong:coso} trình bày chi tiết các cách này.
% EV: J01; J03; J06-J10; C2-08; C2-09; C2-15.

\begin{table}[htbp]
\centering
\caption{Chỗ đứng của khóa luận so với các cách đo và ràng buộc lời hứa. Nội dung viết bằng lời;
“thời gian đi cố định” mô tả mô hình của bài được trích.}
\label{bang:c1-chodung}
\footnotesize
\begin{tabular}{@{}p{2.8cm}p{4.0cm}p{3.3cm}p{4.5cm}@{}}
\toprule
\textbf{Cách tiếp cận} & \textbf{Làm gì; đại lượng xét} &
\textbf{Phần thay đổi khi giữ kế hoạch} & \textbf{Khóa luận kế thừa hay làm khác}\\
\midrule
Fielbaum và Alonso-Mora (2020)~\cite{fielbaum2020}
& Đo thay đổi thời gian; kiểm soát giờ chờ trước quyết định
& Chọn giờ thấp điểm để bỏ qua bất ổn giao thông
& Kế thừa việc đo và kiểm soát; kiểm trần tổng dời giờ trả\\
% EV: J06-J09; C2-06.
Hörl và Zwick (2022)~\cite{horl2022}
& Tính lại tuyến cũ; nới hạn trả theo công thức 7
& Hạn được nới theo giờ cập nhật dưới giao thông hiện tại
& Dùng cùng mốc giữ kế hoạch; trần chung không nới theo giao thông\\
% EV: J01; J03; J05; C2-04.
Schulz và Pfeiffer (2026)~\cite{schulz2026}
& Hạn trả neo vào giờ trong lời mời khách đã nhận
& Thời gian đi được cho trước
& Cài phần hạn chót làm đối chứng, không tái lập toàn bộ phương pháp\\
% EV: J10; J11; C2-07; A09.
Elmi và cộng sự (2026)~\cite{elmi2026}
& Bảo đảm không đổi giờ cho khách được chọn
& Thời gian đi trên mỗi cạnh cố định
& Kế thừa bảo vệ từng khách; giới hạn một tổng qua cả chuyến\\
% EV: C2-08 (nâng J12); A06.
Mt-KaRRi~\cite{laupichler2026}
& Giới hạn thời gian chuyến theo thời gian khách đã chấp nhận
& Thời gian đi trên mỗi cạnh cố định
& Cùng neo vào nghĩa vụ với khách; khác đại lượng bị giới hạn\\
% EV: J32; C2-09; A06.
Lời hứa theo dự báo~\cite{hu2022}
& Hiệu chỉnh giờ đến dự kiến từ dữ liệu
& Tính tới thông tin thời gian thực và lịch sử
& Dùng hồ sơ lịch sử; tính lại lời hứa trên các chuyến đã ghi\\
% EV: J35; C2-15; F01; F03.
Khóa luận
& Trần tổng độ lớn các lần dời do điều phối, cho mỗi khách có lời hứa
& Không tính phần này vào trần; ghi riêng trong sổ
& Đặc tả, cài đặt, kiểm trần và đo đánh đổi với phục vụ\\
% EV: A01; A02; A06; X3-12; C02.
\bottomrule
\end{tabular}
\end{table}

Cách đọc bảng: cùng bảo vệ lời hứa nhưng các cách giới hạn những đại lượng khác nhau. Riêng với Hörl và Zwick,
khóa luận kế thừa mốc tuyến cũ tính lại; phần làm thêm là trần điều phối chung cho mọi khách có lời hứa, cộng
qua cả chuyến và không nới theo tắc đường. Tác dụng thực nghiệm riêng của việc cộng dồn còn chưa được đo.
% EV: J05 (suy luận từ hai định nghĩa); A06; E09.

\section{Mục tiêu và đóng góp}
\label{muc:c1-donggop}

Mục tiêu là đặt một giới hạn mà bộ điều phối kiểm được trước khi chọn kế hoạch, rồi đo giới hạn ấy được gì
và mất gì. Với hãng xe và người đi xe, luật C cho biết mỗi khách có thể bị điều phối dời giờ trả tổng cộng
nhiều nhất bao nhiêu, kể cả khi dự báo sai; đổi lại, bộ điều phối có thể nhận ít yêu cầu hơn.
% EV: A01; A02; A06; A22; C02; I10.

\noindent\textbf{Đóng góp của khóa luận.}
Khóa luận đặc tả và cài đặt trong RideBound luật C: mỗi khách đã được nhận, kể từ khi xác nhận đặt xe,
có một trần cố định cho tổng độ lớn các lần dời giờ trả do điều phối trong cả chuyến; phương án vượt trần bị
loại. Mỗi lần dời được đo so với giờ trả khi giữ kế hoạch cũ, tính lại ở trạng thái hiện tại. Mốc này đã có
ở Hörl và Zwick (2022)~\cite{horl2022}; luật C dùng một trần chung, không nới theo tắc đường và cộng qua cả chuyến.
Luật đi cùng lời hứa theo hồ sơ lịch sử và sổ theo từng khách. Trên 34 kịch bản FleetPy với hệ số tắc đường
toàn thành phố từ dữ liệu Manhattan, cùng 3.672 yêu cầu mỗi cấu hình, trần 30 giây phục vụ ít hơn điều phối
không trần 96 lượt, nhưng thêm 31 lượt không trễ quá một phút so với lời hứa đầu không dự báo.
Luật hạn chót 30 giây phục vụ cùng số lượt, nhưng tính cả phần thay đổi do giao thông, nên phải phá luật của
mình với ít nhất một khách ở 28 trong 34 kịch bản. Luật C không tính phần đó vào trần. Tính không bị ép của
luật C cần giả thiết (N) ở mục~\ref{muc:giunguyen} và phần trần đã dùng không vượt giới hạn; điều phối không
trần cũng có 0 trong 20.760 quyết định bị ép. Luật C không giữ lời hứa đầu tốt hơn: đo cùng ngưỡng trễ
30 giây, luật C có 292 trong 2.338 khách được phục vụ (12,5\%), còn luật hạn chót có 289 trong 2.338 (12,4\%),
gộp theo khách. Khi có dự báo, đệm thêm 30 giây cho điều phối không trần cho cùng số lượt không trễ quá
một phút như luật C và nhiều hơn ở ngưỡng hai phút; lời hứa dài thêm trung bình 37,2 thay vì 7,4 giây,
gộp trên 2.434 và 2.338 khách được phục vụ. Trong 132 lần chạy luật C đã kiểm toàn bộ bản ghi chặn,
3 trong 1.780 bản ghi chỉ bị chặn do cộng dồn, đều ở trần 120 giây; tác dụng lên kết quả chưa được đo.
% EV: A06; A13; X3-20; A01; A02; J01; J03; J05; F01; X3-12; B06; B09; C01; C02; C08; C04; N-21; A22; X3-27; X3-16; C4A-01; C10; N-16; N-18; N-15; N-08; E09.

Các kết quả chính dưới đây là số đếm gộp, không phải trung bình theo kịch bản. Mọi số lượt phục vụ hay không trễ
đều xét trên cùng 3.672 yêu cầu; “không trễ” gồm cả khách tới sớm.
% EV: C01; C06; C08; N-16; N-06; C4A-01.
\begin{itemize}[leftmargin=0.7cm,itemsep=2pt,topsep=3pt]
\item Phục vụ: không trần 2.434 lượt; luật C và luật hạn chót, cùng trần 30 giây, đều 2.338 lượt, có 2.332 lượt chung.
% EV: C01; C04.
\item Không trễ quá một phút, chưa dự báo: không trần 2.231 lượt, luật C 2.262; khách trễ là 203 trong 2.434 và
76 trong 2.338; khách đi chung là 555 trong 2.434 và 313 trong 2.338.
% EV: C06; C08; N-09.
\item Có dự báo: luật C và không trần đệm 30 giây cùng 2.328 lượt không trễ quá một phút; ở hai phút là 2.338
và 2.366. Tới sớm quá một phút là 4 trong 2.338 và 44 trong 2.434 khách.
% EV: C10; N-16; N4B-09.
\item Trần 120 giây: phục vụ 2.387 lượt; không trễ quá một và hai phút là 2.255 và 2.366; có dự báo là 2.322 và
2.379, với lời hứa dài thêm trung bình 7,3 giây trên 2.387 khách.
% EV: N-06; N-16.
\item Cộng dồn: tập 132 lần chạy đã kiểm gồm 98 lần cũ và 34 lần trần 120 giây; chưa chạy đối chứng mỗi lần một trần.
% EV: N-15; N-08; E09.
\end{itemize}

\begin{khungy}
\textbf{Khóa luận không nhận là mới.}
Chèn khách động, giới hạn giờ đến, gán lại xe, tương đồng tuyến, cam kết tối thiểu, nhất quán thời gian, đo và
kiểm soát độ bất ổn, bảo vệ từng khách và lời hứa theo dự báo đều có tiền lệ. Phần khóa luận làm là đặc tả,
cài đặt và kiểm trần đã nêu, cùng sổ và phép đo đánh đổi. Luật C không bảo đảm giờ đến thực tế; các kết quả
không xác nhận lợi ích kinh doanh hay mức hài lòng của khách.
\end{khungy}
% EV: J38; J01; J03; J06-J10; C2-08; C2-15; A06; X3-12; K05; I07; I10; K38.

\section{Phạm vi}
\label{muc:c1-phamvi}

Phạm vi đánh giá được giới hạn như sau:
\begin{itemize}[leftmargin=0.7cm,itemsep=3pt,topsep=3pt]
\item \textbf{Môi trường:} RideBound nối với FleetPy 1.0.2 và bộ dữ liệu chuẩn Manhattan.
Đồ thị 96 nút có sẵn thời gian đi cho từng cặp điểm, không phải mạng đường phố đầy đủ.
% EV: B01; B02; B03; K48.
\item \textbf{Quy mô:} mỗi kịch bản có 108 yêu cầu, 8 xe 4 chỗ; đón trong vòng 600 giây kể từ lúc gọi,
ngồi xe không quá 1,5 lần thời gian đi thẳng. Đây là cấu hình thử, không phải chuẩn ngành.
% EV: B02; B03.
\item \textbf{Tắc đường:} một hệ số chung cho cả thành phố, đổi theo khung 15 phút; tốc độ nền lấy lúc 00:00.
Mô hình này không mô tả tắc đường riêng trên từng đoạn.
% EV: B06; B07; K49.
\item \textbf{Dữ liệu kiểm tra:} tập chính gồm 34 kịch bản của năm ngày 14--18/11/2018.
Các khung giờ và mẫu nhu cầu có phần chồng nhau, nên không phải 34 quan sát độc lập.
% EV: B08; B09; B11; I05.
\item \textbf{Ngoài phạm vi:} không có mô hình khách hủy chuyến hay từ chối lời mời; không kiểm khả năng mở rộng
sang mạng lớn hoặc hiệu quả kinh doanh. Dự báo chính chỉ được tính lại trên các chuyến đã ghi, không đổi quyết định điều phối.
% EV: I07; B04; I01; I10; F03; X3-12.
\end{itemize}

\section{Cấu trúc báo cáo}
\label{muc:c1-cautruc}

Chương~\ref{chuong:coso} trình bày bài toán và các nghiên cứu liên quan.
Chương~\ref{chuong:phuongphap} nêu luật C, lời hứa theo dự báo và sổ theo từng khách.
Chương~\ref{chuong:thucnghiem} đo đánh đổi, đối chiếu các luật và nêu hạn chế.
Chương~\ref{chuong:ketluan} tổng hợp kết quả và các hướng kiểm tra tiếp.
% EV: — (lộ trình báo cáo; nhãn cuối theo brief ROUND31).
